\documentclass{article}
% formatting
\usepackage{aux/style_files/arxiv}
\usepackage[utf8]{inputenc} % allow utf-8 input
\usepackage[T1]{fontenc}    % use 8-bit T1 fonts
% cross-referencing
\usepackage{hyperref}
\usepackage{url} 
% tables
\usepackage{booktabs}
\usepackage{tabularx}
\usepackage{caption} 
\captionsetup[table]{skip=10pt}
% fonts
\usepackage{amsfonts}
\usepackage{microtype}
% figures
\usepackage{graphicx}
% enumeration
\usepackage{enumitem}



\title{Life-Cycle Assessment for Sustainable Aviation}


\author{
    \href{https://orcid.org/0000-0003-4859-2650}
    {\includegraphics[scale=0.06]{aux/figures/orcid.pdf}
    \hspace{1mm}
    Michael Weinold} \\
	Paul Scherrer Institut\\
	Laboratory for Energy Systems Analysis\\
	Group for Technology Assessment\\
	OSHA/D22\\
    5232 Villigen PSI \\
    Switzerland \\
	\texttt{michael.weinold@psi.ch} \\
}

\date{28. July 2022}

\renewcommand{\headeright}{Doctoral Research Plan}
\renewcommand{\undertitle}{Doctoral Research Plan}

%%% Add PDF metadata to help others organize their library
%%% Once the PDF is generated, you can check the metadata with
%%% $ pdfinfo template.pdf
\hypersetup{
pdftitle={A template for the arxiv style},
%pdfsubject={q-bio.NC, q-bio.QM},
pdfauthor={Michael Weinold},
%pdfkeywords={First keyword, Second keyword, More},
}

\begin{document}
\maketitle

\begin{abstract}
	\cite{becattini_role_2021}
\end{abstract}


\section{Introduction}

\cite{prussi_corsia_2021}

how are sustainable aviation fuels defined -> CORSIA \cite{prussi_corsia_2021}

    Carbon emissions estimations for individual flights have been available from major flight search providers, including \textit{Google Flights} \cite{holden_google_2021}, \textit{Kayak} \cite{noauthor_kayak_2021} and \textit{SkyScanner} \cite{crosthwaite_how_2021} from 2021. Early studies have shown that the prominent display of carbon emissions can have positive enironmental impact by nudging consumers to choose flights that XXX \cite{amenta_adding_2020}\cite{sanguinetti_nudging_2022}.
    
    However, different methodologies are used to estimate carbon emissions, leading to significant differences in total estimated carbon emissions. For instance, the development team of \textit{Google Flights} changed the aclculation methodology to account for direct-air carbon emissions only, disregarding non-carbon emissions \cite{ali_commit_2022}. This has drawn criticism of \textit{"airbrushing emissions"} by environmental non-governmental organizations \cite{hern_google_2022}\cite{rowlatt_google_2022}.

\section{Hybrid Life Cycle Assessment}

	What is life-cycle assessment?
	Why is its use in aviation a novelty?
	
	Some recent publications have completely disregarded the impact of life-cycle emissions under the assumption that these would be negligible \cite{brazzola_definitions_2022}.
	
	Why is the use of hybrid life-cycle inventory methods in aviation a novelty?
	
\section{USP}
	
    Existing knowledge within PSI, for instance the recent thesis by Saad \cite{saad_synthetic_2022} and unpublished work by \cite{sacchi_climate-neutral_2022}.

\section{Ideas}

    Connect this to "M2.1: Publication analysing impact of multiple standards and standards implementations in LCA"

    Connect this to "RECCE 2035: Resource to Climate Comparison Evaluator"

	Chris: research and development phase often missed completely in LCA
	Chris: Airbus is using Brightway in future aircraft design
	Chris: knowledge in the group on sustainable fuels (Tom Terlouw?)
	Laura: Cambridge Aviation Impact Accelerator meeting in September - maybe meet before that?

\section{Questions}

How do the results from Romain's preprint \cite{sacchi_climate-neutral_2022} compare against the new ICAO paper \cite{prussi_corsia_2021}

\newpage
\section{Systematic Literature Review}

    \subsection{Life-Cycle Assessment of (Sustainable) Aviation}

\newpage
\section{Appendix}

 % https://docs.google.com/spreadsheets/d/14gmD96CXOYMKjaUK2KIvosC4LcfpzAc-_INdCU4LAro/edit#gid=1019504979

    \subsection{Details of Systematic Literature Review Methodology}
    
        This work uses the \textit{SALSA} (\textbf{S}earch, \textbf{A}ppraisa\textbf{L}, \textbf{S}ynthesis, \textbf{A}nalysis) framework rationalization introduced by Grant et al. \cite{grant_typology_2009}. Search was 
    
        \subsubsection{Boolean Search Logic}

            \begin{table}[htbp]
                \centering
                \begin{tabularx}{\textwidth}{| X | X |}
                    \hline
                    \textit{field of research} & \textit{search logic} \\
                    \hline
                    (1.1) LCA specific to (sustainable) aviation & \texttt{("life cycle assessment" OR "life-cycle assessment") AND ("aviation" OR "sustainable aviation")} \\
                    \hline
                    (1.2) hybrid LCA specific to (sustainable) aviation &  \texttt{("life cycle assessment" OR "life-cycle assessment") AND ("aviation" OR "sustainable aviation")} \\
                    \hline
                    (2.1) hybrid LCA methods & \texttt{("life cycle assessment" OR "life-cycle assessment") AND ("aviation" OR "sustainable aviation")} \\
                    \hline
                    (2.2) prospective (hybrid) LCA methods & \texttt{("life cycle assessment" OR "life-cycle assessment") AND ("aviation" OR "sustainable aviation")} \\
                    \hline
                    (3.1) impact of different GHG standards/implementations in LCA & \texttt{("life cycle assessment" OR "life-cycle assessment") AND ("aviation" OR "sustainable aviation")} \\
                    \hline
                \end{tabularx}
                \caption{Boolean search logic for the different fields of research considered in this work. For the inclusion and exclusion considerations that informed the logic, compare table \ref{tab_inclusion_exclusion}.}
                \label{tab_bool}
            \end{table}
            
        \subsubsection{Inclusion/Exclusion Criteria}
        
            \begin{enumerate}
                \item LCA specific to (sustainable) aviation
                \begin{enumerate}[label*=\arabic*]
                    \item LCA used as primary method of investigation
                    \item LCA XXX no catering etc.
                \end{enumerate}
            \end{enumerate}
        
            \begin{table}[htbp]
                \centering
                \begin{tabularx}{\textwidth}{| X | X |}
                    \hline
                \multicolumn{2}{|l|}{(1.1) LCA specific to (sustainable) aviation}  \\
                    \hline
                    \textit{inclusion criteria} & \textit{exclusion criteria} \\
                    \hline
                        LCA used as primary method  \newline
                    &
                        LCA target domain too specific \newline (e.g. "aviation catering") \newline
                        fuel input too specific \newline (e.g. "bioethanol production from expired cookies") \newline
                        journal subject ares in social sciences \newline
                        journal subject area in life sciences
                    \\
                    \hline
                \multicolumn{2}{|l|}{(1.2) hybrid LCA specific to (sustainable) aviation}  \\
                    \hline
                    \textit{inclusion criteria} & \textit{exclusion criteria} \\
                    \hline
                        hybrid LCA used as primary method  \newline
                    &
                        all exclusion criteria applied in (1.1) \newline
                        primary focus on hybrid-electric propulsion \newline
                        (e.g. "Feasibility Study of Hybrid Propulsion")
                    \\
                    \hline
                \multicolumn{2}{|l|}{(2.1) hybrid LCA methods}  \\
                    \hline
                    \textit{inclusion criteria} & \textit{exclusion criteria} \\
                    \hline
                        hybridization methodology is primary focus \newline
                        process-based life-cycle assessment (PLCA) \newline
                        multi-regional input-output tables (MRIO) \newline
                        environmentally extended IO analysis (EEIOA) \newline
                        specific databases: EXIOBASE, ecoinvent \newline
                        focus on truncation errors and double-counting
                    &
                        journal subject area in health sciences \newline
                        journal subject area in social sciences \newline
                        journal subject area in life sciences
                    \\
                    \hline
                \multicolumn{2}{|l|}{(2.2) hybrid LCA methods}  \\
                    \hline
                    \textit{inclusion criteria} & \textit{exclusion criteria} \\
                    \hline
                        all inclusion criteria applied in (2.1) \newline
                        prospective/future hybrid LCA
                    &
                        all exclusion criteria applied in (2.1)
                    \\
                    \hline
                \multicolumn{2}{|l|}{(3.1) impact of different GHG standards/implementations in LCA}  \\
                    \hline
                    \textit{inclusion criteria} & \textit{exclusion criteria} \\
                    \hline
                        comparison of specific assessment frameworks in the context of LCA\newline
                        (i.e. ISO 14040-14044, GHG Protocol, ILCD Handbook, ISO 14067, CDP, GRI, SBTi, PEF/OEF, ESG)
                    &
                        narrow focus on any one framework from the inclusion criteria \newline
                        LCA not in research context
                    \\
                    \hline
                \end{tabularx}
                \caption{Inclusion and exclusion criteria for systematic literature review. Where applicable and possible, these were translated to a boolean search logic. Compare also table \ref{tab_bool}.}
                \label{tab_inclusion_exclusion}
            \end{table}

\newpage

\bibliographystyle{IEEEtran}
\bibliography{references}

\end{document}
